\documentclass[12pt,a4paper,pdftex,final,oneside,titlepage]{book}
\usepackage[pdftex]{graphicx,color}
\usepackage{graphicx}
\usepackage[left=4cm,top=3cm,right=3cm,bottom=3cm,includeheadfoot]{geometry} %untuk mengatur margin
\usepackage[pdftex]{hyperref}
\usepackage{lscape}
%\begin{landscape} untuk awal landsacpe
%landscape content
%\end{landscape} untuk akhir landscape
\usepackage{longtable}
\usepackage{setspace} %untuk mengatur spasi, pilihan: \doublespacing \singlespacing atau \onehalfspacing   
\usepackage[bahasa]{babel}
\selectlanguage{bahasa}
\usepackage{fancyhdr,lastpage}
%\usepackage[Lenny]{fncychap}
%\usepackage[Sonny]{fncychap}
\usepackage{pstricks}
\usepackage{listings} %listing program
\usepackage{appendix}
\usepackage{subfig} %sub figure
\usepackage{sidecap} %side caption
\usepackage{ulem} %striketrough
%==========================================================================================================
% Bagian input data untuk keterangan file PDF yang dihasilkan
%========================================================================================================== 
\pdfinfo{
/Author (Rizky Zulkarnaen)
/CreationDate (D:20100602300000) % D:YYYYMMDDhhmmss (diisi dengan tanggal pembuatan skripsi)
/ModDate (D:20100602300000) % diisi dengan tanggal modifikasi skripsi
/Creator (LaTeX)
/Producer (pdfLaTeX)
/Title (TA_AR) % diisi dengan nama file skripsi
/Subject (Web-based Interactive 3D Model Viewer With Augmented Reality) % diisi dengan judul skripsi
/Keywords (augmented reality, computer vision, 3D) % diisi dengan kata-kata kunci isi skripsi
}

%==========================================================================================================
% Bagian pengaturan hyperref
%========================================================================================================== 
\hypersetup{
pdftitle={Web-based Interactive 3D Model Viewer With Augmented Reality}, % diisi dengan judul akhir skripsi
pdfauthor={Rizky Zulkarnaen}, % diisi dengan nama penulis skripsi
pdfkeywords={augmented reality computer vision 3D}, % diisi dengan kata-kata kunci dari isi skripsi
bookmarksnumbered,
pdfstartview={FitH},
	citecolor = black,
	linkcolor=black,
  colorlinks
}

%==========================================================================================================
% Koleksi pemenggalan kata, isi bagian ini bisa ditambah jika Anda mendapati kata yang salah pemenggalannya 
%==========================================================================================================
\hyphenation{client}

%==========================================================================================================
% Bagian pengaturan paragraf 
%==========================================================================================================
\setlength{\parindent}{0in} % jika awal paragraf ingin dibuat menjorok ke dalam, ganti angka 0 nya. 
\setlength{\parskip}{6pt plus1pt minus1pt}
\newcommand{\marginal}[1]{\leavevmode\marginpar{\color{cyan}\footnotesize\raggedright#1\par}} 
\onehalfspacing
%==========================================================================================================
% Bagian pengaturan Fancy Header
%==========================================================================================================
\setlength{\headheight}{15pt}

%==========================================================================================================
% BAGIAN ISI DOKUMEN 
%==========================================================================================================
\begin{document}  

%==========================================================================================================
% Bagian pembuatan halaman judul 
%==========================================================================================================
\pagestyle{empty} 
	\begin{titlepage} 
		\begin{center} 
		{\Large\textbf{Tugas Akhir}}\\
		\par
		\vspace{2cm}
		\textbf{\Large Web-based Interactive 3D Model Viewer With Augmented Reality} % diisi dengan judul skripsi
		\vspace{0.5cm} \\
		oleh :\\[0.4cm] 
		{\large\underline{Rizky Zulkarnaen}}\\% ganti Nama Penulis dengan nama Anda sendiri
		$050402056$ \\% ganti NIM
		\end{center}
		
		\par
		\vfill
		
		\begin{figure}[h]
			\begin{center}
				\includegraphics[width=4.2cm]{./images/logoUSU} % isi dengan nama file logo
			\end{center}
        \end{figure}
        
        \vfill 
        
      \begin{center}
      {\bf \large DEPARTEMEN TEKNIK ELEKTRO}\\[0.4cm] 
      {\bf \large FAKULTAS TEKNIK}\\[0.4cm] 
      {\bf \large UNIVERSITAS SUMATERA UTARA}\\[0.4cm] 
      {\bf \large MEDAN}\\[0.4cm] 
      {\bf \large 2010}
      \end{center}
      
   \end{titlepage} 

%==========================================================================================================
% Bagian pembuatan halaman pengesahan skripsi 
%==========================================================================================================
\thispagestyle{empty} 
\begin{center}{%\fontfamily{phv}\selectfont

\textbf{\large Web-based Interactive 3D Model Viewer With Augmented Reality}\\[1cm] 

\par 
\vspace{0.4cm} 
oleh : 
\\[1cm] 
{\large\underline{Rizky Zulkarnaen}}\\ 
$050402056$
} 
\vspace{1.4cm} 

This final task proposed to complete requirement bachelor degree in technique \\[0.2cm] 

Agreed by:\\ 
Pembimbing\\[2.2cm] 
\textbf{\bfseries \underline{Ori Novanda, ST}}\\[0.2cm] 
NIP:\\[0.4cm]
Known by:\\ 
Pelaksana Harian Ketua Departemen Teknik Elektro\\[2.2cm] 
\textbf{\bfseries \underline{Prof. Dr. Ir. Usman S. Baafai}}\\[0.2cm] 
NIP:19461022197302001\\ 

\end{center} 
%========================================================================================================== 
% Bagian pengaturan header dan footer
%========================================================================================================== 
\pagestyle{fancy}
\lhead{} \chead{} \rhead{\thepage}
\lfoot{} \cfoot{} \rfoot{}
\renewcommand\headrulewidth{0pt}
\renewcommand\footrulewidth{0pt}

%==========================================================================================================
% ISI SKRIPSI 
%==========================================================================================================
\pagenumbering{roman}
\setcounter{tocdepth}{5} 
\setcounter{secnumdepth}{3} %ditambahkan agar perintah \subsubsection memiliki angka

%==========================================================================================================
%ABSTRAK
%==========================================================================================================

\chapter*{\centering Abstract}
\thispagestyle{empty}
\addcontentsline{toc}{chapter}{Abstrak} 
Today, it has become very popular to represent 3D CG models onto the screen of a workstation or personal computer. In almost all software that displays 3D models, users can control their viewpoint of projection with pointing devices such as a mouse or trackball. However it is not intuitive to control the viewpoint with such devices, because there is an insufficient degree of freedom. Although locating viewpoint needs 6 degrees of freedom (3 degrees in position, and 3 degrees in direction), off-the-shelf pointing devices are basically two-dimensional controllers. We propose a 3D model viewer with a real-time marker tracking system. This system provides an easy method to make observations of 3D models from any direction and to any scale that users may choose. 

\vspace*{2cm}

\textbf{Kata Kunci : augmented reality, computer vision, image recognition}

%==========================================================================================================
% AWAL KATA PENGANTAR  
%==========================================================================================================
\chapter*{KATA PENGANTAR}
\addcontentsline{toc}{chapter}{Kata Pengantar} 

Puji dan syukur kepada Allah SWT penulis ucapkan, karena hanya dengan kehendakNya penulis dapat menyelesaikan Tugas Akhir ini, yang berjudul "Perancangan Aplikasi Viewer Model 3D Interaktif Berbasis Web dengan Teknologi Augmented Reality". Tak lupa pula shalawat dan salam kepada Junjungan kita Nabi Besar Muhammad SAW yang telah membimbing kita kepada Islam. Di samping itu, penulis mengucapkan terima kasih yang sebesar-besarnya kepada:

\begin{enumerate}
\item Ayahanda Zulkifli. Y dan Ibunda Nurmawati yang telah membesarkan dan memberikan kasih sayang yang tulus kepada penulis.
\item 2.
\item Bapak Ori Novanda ST, MT, selaku dosen pembimbing penulis yang telah sangat banyak membantu dalam penulisan Tugas Akhir ini.
\item Bapak Ir. Riswan Dinzi selaku dosen wali penulis, atas bimbingannya selama penulis kuliah di Departemen Teknik Elektro, Fakultas Teknik, Universitas Sumatera Utara.
\item Bapak Prof. Dr. Ir. Usman S. Baafai dan Bapak Rahmad Fauzi, ST, MT, selaku Pelaksana Harian Ketua dan Sekretaris Departemen Teknik Elektro, Fakultas Teknik, Universitas Sumatera Utara.
\item Seluruh staf pengajar dan pegawai Departemen Teknik Elektro Fakultas Teknik Universitas Sumatera Utara.
\item Asisten Laboratorium Sistem Pengaturan dan Komputer, Salman Alfarisi dan Fachrurozi Nasution, yang menemani dan banyak membantu penulis dalam pengerjaan Tugas Akhir ini.
\item Teman-teman stambuk 2005,  dan teman-teman yang belum disebut namanya, yang selama ini menjadi teman diskusi di kampus.

\end{enumerate}
Berbagai usaha telah penulis lakukan demi terselesaikannya Tugas Akhir ini dengan baik, tetapi penulis menyadari akan kekurangan dan keterbatasan penulis. Oleh karena itu, saran dan kritik dengan tujuan menyempurnakan dan mengembangkan kajian dalam bidang ini sangat penulis harapkan.

Akhir kata penulis berharap agar Tugas Akhir ini dapat bermanfaat bagi pembaca dan penulis.

\vspace*{2cm} 
\begin{flushright} 
Medan,\hspace{0.2cm} Juni 2010 \hspace{1cm} \\ % sesuaikan dengan tempat dan tanggal penulisan skripsi 
Penulis,\hspace{2.2cm}
\par 
\vspace*{1cm} 
\par \underline{Rizky Zulkarnaen}\\%\hspace{0.5cm}
NIM:$050402056$ \hspace{0.6cm} \par
\end{flushright} 
%==========================================================================================================
% AKHIR KATA PENGANTAR 
%==========================================================================================================

\clearpage
\addcontentsline{toc}{chapter}{List of Contents} 
\tableofcontents  %perintah untuk menampilkan daftar isi

\clearpage
\addcontentsline{toc}{chapter}{List of Tables} 
\listoftables %perintah untuk menampilkan daftar tabel 

\clearpage
\addcontentsline{toc}{chapter}{List of Figures} 
\listoffigures %perintah untuk menampilkan daftar gambar  
%==========================================================================================================
% MULAI BAB I
%==========================================================================================================
\chapter{INTRODUCTION} %JUDUL BAB 
\label{chap:pendahuluan} 
\pagenumbering{arabic} 

\section{Background} %Judul subbab 
\label{sec:latar_belakang} 
%isi subbab latar belakang
Modelling 3 dimension object (3 D) very needed in every application, whether for simulation or model introduction from real object which difficult to presence in physics, because of time and space limit. The model of real object able to presence in virtual which can be display through screen or computer, thus modeling an object easy to do with low cost. Many of fields need modeling of 3D virtual object, such as body organ modeling helpfull in medical, building modeling, modeling product which to sell, and so on.

Modeling virtual object need good interaction with user to earn clear viewpoint from the object. Some of 3D modeling use pointer, such as keyboard or mouse that still lack in give degree of freedom viewpoint from the 3D model. Augmented reality is a merger concept between virtual reality and world reality. There by, 2 dimension virtual objects (2D) or 3 dimension (3D) like real and unite to the real world. In augmented reality technology, user able to see real world surround with virtual object added that comes from computer. With augmented reality technology, 3D model presentation able to presence more interactive. Object in the real world can be used as tracker for interaction needed in 3D virtual model, so helpful in appearance that virtual object.

Computer application (software) which can show 3D model, usually need installation process in order to operated, and sometimes only operate in specific platform. The development of web technology enable many computer's applications make in web based. Applicatons with web based is very easy to access, moreover it can be operate in various platform. Application with difficult specification through specific web plugin can be made and can be access only from web browser. Requirement of web based application will ever increasing, because of user's application like to use easy application and without need of installation.  Likewise, with application for viewer model 3D, web based application viewer model 3D will facilitate people who want to look 3D model in an interactive manner through browser without need of application installation or may be for companies, which want to show their product in 3D model as interactively to the consumer.

\section{Writing Purpose}
\label{sec:tujuan_penulisan}
The aim of this research is to design an application that facilitate object 3D presentation (viewer model 3D), especially interaction with applying augmented reality technology. Application that will be made is web based with Rich Internet Application (RIA) concept, in this case is used technology from Adobe Flex that have integration to the web, so the application can be access through web browser and certainly can be access in various platform. This application also use database to save 3D object and show it again.

\section{Tinjauan Pustaka}
\label{sec:tinjauan_pustaka}
tinjauan pustaka
%==============================================================================================
% Subbab Rumusan Masalah
%==============================================================================================
\section{Formulation of Problem} %Judul subbab
\label{rumusan}
from the background that have explain above, can be defined some problem, that is:\\
\begin{enumerate}
\item How augmented reality technology applied to make interactive viewer model 3D.
\item Design web based application that able to save 3D object into a database and then show again with Rich Internet Application (RIA) concept.
\end{enumerate}

%==============================================================================================
% Subbab Batasan Masalah
%==============================================================================================
\section{Limit of Problem} %Judul subbab
\label{batasan}
To avoid broader investigation in this study, it has limited in some cases:\\
\begin{enumerate}
\item Study augmented reality concept in summary and supporting theory in design viewer 3D application with augmented reality.
\item Not study the making of 3D object, but only study how to show the 3D object with augmented reality support.
\item The making of the applicaton is web based with utilize Adobe Flex technology. That application able to save 3D object that being input by user, determine marker for 3D object, then show the 3D object with augmented reality technology.
\item Not study detail database system that use to save 3D data object.
\end{enumerate}

\section{Research Methodology}
\label{sec:metodologi_penelitian}
Research Methodology

\section{Writing Systematic}
\label{sec:sistematika_penulisan}
The writing of this final task, presence with writing systematic:
\par\noindent
\begin{longtable}{llp{11cm}}
CHAPTER I & : & INTRODUCTION\\
& & This chapter present view of all matter that studied in this final task, study about background, research aim, limit of problem, research method, and writing systematic.\\
CHAPTER II & : & DASAR TEORI\\
& & This chapter study augmented reality, supporting theories of augmented reality especially augmented reality with use marker and also Rich Internet Application concept that use Adobe Flex.\\
CHAPTER III & : & APPLICATION DESIGN AND IMPLEMENTATION\\
& & This chapter study about the design application, the use of Rich Internet Application (RIA) technology that is Adobe Flex and also library for augmented reality.\\
CHAPTER IV & : & TESTING AND ANALYSIS\\
& & This chapter study about testing and analysis of the problem.\\

CHAPTER V & : & CONCLUSION AND SUGGESTION\\
& & This chapter content conclusion and suggestion from the application that have designed.

\end{longtable}
%==========================================================================================================
%Penulisan BAB di file yang berbeda
%==========================================================================================================
%\chapter{PENDAHULUAN} %JUDUL BAB 
\label{chap:pendahuluan} 
\pagenumbering{arabic} 

\section{Latar Belakang} %Judul subbab 
\label{sec:latar_belakang} 
%isi subbab latar belakang
Pemodelan objek tiga dimensi (3D) sangat diperlukan dalam berbagai aplikasi, baik untuk simulasi maupun untuk pengenalan model dari objek nyata yang sulit disajikan secara fisik dikarenakan keterbatasan ruang dan waktu. Model suatu objek nyata dapat disajikan secara virtual yang dapat dilihat melalui suatu layar atau \textit{display} dengan bantuan komputer sehingga pemodelan suatu objek mudah dilakukan dengan biaya yang murah. Banyak bidang yang memerlukan pemodelan objek virtual 3D ini, misalnya pemodelan organ tubuh yang bermanfaat dalam dunia kedokteran, pemodelan bangunan, pemodelan suatu produk yang akan dijual, dan lain sebagainya.

Pemodelan objek virtual tersebut memerlukan interaksi yang baik dengan pengguna untuk mendapatkan sudut pandang yang jelas dari objek tersebut. Beberapa \textit{software} pemodelan 3D menggunakan \textit{pointer} seperti \textit{keyboard} ataupun \textit{mouse} yang dirasakan masih kurang dalam memberikan derajat kebebasan sudut pandang dari model 3D tersebut. \textit{Augmented Reality} merupakan suatu konsep perpaduan antara \textit{"virtual reality"} dengan \textit{"world reality"}. Sehingga objek-objek virtual 2 Dimensi (2D) atau 3 Dimensi (3D) seolah-olah terlihat nyata dan menyatu dengan dunia nyata. Pada teknologi \textit{Augmented Reality}, pengguna dapat melihat dunia nyata yang ada di sekelilingnya dengan penambahan obyek virtual yang dihasilkan oleh komputer.  Dengan teknologi \textit{Augmented Reality}, penyajian model 3D dapat disajikan dengan lebih interaktif. Objek di dunia nyata dapat digunakan sebagai \textit{tracker} untuk keperluan interaksi model 3D virtual sehingga memudahkan dalam menampilkan objek virtual tersebut.

Aplikasi komputer (\textit{software}) yang dapat menampilkan model 3D \textit{viewer model 3D} biasanya memerlukan proses instalasi agar dapat dijalankan, dan terkadang hanya dapat berjalan di platform tertentu saja. Dengan berkembangnya teknologi \textit{web}, banyak aplikasi komputer yang dibuat berbasis \textit{web}. Aplikasi yang berbasis \textit{web} sangat mudah diakses, bahkan dapat berjalan di berbagai platform. Dengan \textit{web plug-in} tertentu, aplikasi yang tergolong rumit pun bisa dibuat dan dapat diakses hanya lewat \textit{web browser}. Kebutuhan akan aplikasi yang berbasis \textit{web} ini akan terus meningkat, karena pengguna menyukai aplikasi yang mudah dijalankan dan tanpa perlu instalasi. Begitu pula dengan aplikasi untuk \textit{viewer} model 3D, aplikasi \textit{viewer} 3D berbasis \textit{web} akan memudahkan seseorang yang ingin melihat model 3D secara interaktif melalui \textit{web browser} tanpa perlu instalasi aplikasi atau mungkin bagi perusahaan yang ingin menampilkan produknya dalam bentuk model 3D secara interaktif kepada konsumen.

Berdasarkan uraian tersebut, penulis mengembangkan suatu aplikasi untuk menampilkan objek 3D virtual dan dapat diakses dengan mudah dari mana saja menggunakan \textit{web browser} tanpa perlu instalasi aplikasinya. Pembahasan pengembangan aplikasi ini dibuat menjadi tugas akhir yang diberi judul \textbf{``Perancangan Aplikasi \textit{Viewer} Model 3D Interaktif Berbasis \textit{Web} dengan Teknologi \textit{Augmented Reality}"}.

\section{Tujuan Penulisan}
\label{sec:tujuan_penulisan}
Tujuan dari penulisan tugas akhir ini adalah untuk merancang suatu aplikasi yang memudahkan penyajian objek 3D (\textit{viewer} model 3D) terutama interaksinya dengan menerapkan teknologi \textit{Augmented Reality}. Aplikasi yang akan dibuat ini berbasis \textit{web} dengan konsep \textit{Rich Internet Application (RIAs)}, sehingga aplikasi tersebut dapat menampilkan objek 3D virtual melalui \textit{web browser} dan tentunya dapat diakses di  berbagai platform tanpa perlu instalasi aplikasi di sistem pengguna. 

%==============================================================================================
% Subbab Rumusan Masalah
%==============================================================================================
\section{Rumusan Masalah} %Judul subbab
\label{rumusan}
Dari latar belakang masalah, maka dapat dirumuskan beberapa permasalahan, yaitu sebagai berikut:
\begin{enumerate}
\item Menerapkan teknologi \textit{augmented reality} untuk membuat \textit{viewer} model 3D yang interaktif.
\item Merancang aplikasi berbasis \textit{web} untuk \textit{viewer} objek 3D dengan konsep \textit{Rich Internet Application} (RIAs) dan teknologi \textit{Augmented Reality}.
\end{enumerate}

%==============================================================================================
% Subbab Batasan Masalah
%==============================================================================================
\section{Batasan Masalah} %Judul subbab
\label{batasan}
Untuk menghindari pembahasan yang terlalu luas, maka penulis akan membatasi tugas akhir ini dengan hal-hal sebagai berikut:
\begin{enumerate}
\item Membahas secara ringkas konsep \textit{augmented reality} dan teori  yang menunjang perancangan aplikasi viewer 3D dengan \textit{augmented reality}.
\item Tidak membahas pembuatan objek 3D nya, tetapi hanya membahas bagaimana menampilkan objek 3D tersebut dengan bantuan \textit{engine/library} 3D dan teknologi augmented reality.
\item Pembuatan aplikasinya berbasis \textit{web} yang dapat diakses dari \textit{web browser}. Aplikasi tersebut nantinya dapat menampilkan objek 3D dan bisa berinteraksi dengan pengguna menggunakan sebuah \textit{tracker}.
\end{enumerate}

\pagebreak

\section{Metodologi Penulisan}
\label{sec:metodologi_penulisan}
Penulisan dilakukan dalam beberapa tahapan:
\begin{enumerate}
\item Tahap penelitian platform dan \textit{framework}\\
Pada tahap ini penulis meneliti platform yang akan digunakan, serta juga \textit{framework}/\textit{library} yang cocok untuk pengembangan aplikasi tersebut.
\item Tahap perancangan aplikasi\\
Pada tahap ini dirancang aplikasinya dengan menggunakan \textit{framework} yang telah ditentukan di tahap penelitian. 
\item Tahap pengembangan aplikasi\\
Pada tahap ini penulis membangun aplikasinya sesuai dengan rancangan.
\item Tahap implementasi dan pengujian\\
Pada tahap ini penulis mengimplementasikan aplikasi dan melakukan pengujian dari aplikasi yang telah dikembangkan.
\end{enumerate}

\section{Sistematika Penulisan}
\label{sec:sistematika_penulisan}
Penulisan Tugas Akhir ini disajikan dengan sistematika penulisan sebagai berikut:
\par\noindent
\begin{longtable}{llp{11cm}}
BAB I & : & PENDAHULUAN\\
& & Bab ini merupakan gambaran menyeluruh tentang apa yang diuraikan dalam Tugas Akhir ini, yaitu pembahasan tentang latar belakang, tujuan penulisan, batasan masalah, metode penulisan, dan sistematika penulisan.\\
BAB II & : & DASAR TEORI\\
& & Bab ini membahas tentang \textit{augmented reality}, teori-teori yang mendukung \textit{augmented reality}, teknik \textit{display} \textit{augmented reality}, konsep aplikasi berbasis \textit{web} dan juga library untuk \textit{augmented reality}.\\
BAB III & : & PERANCANGAN DAN IMPLEMENTASI\\
\nopagebreak
& & Bab ini merupakan bab yang membahas perancangan aplikasinya, dimulai dari gambaran umum aplikasi desain aplikasi dan terakhir implementasi aplikasi.\\
BAB IV & : & PENGUJIAN DAN ANALISA\\
\nopagebreak
& & Bab ini membahas tentang pengujian dan analisa dari aplikasi yang dibangun.\\
BAB V & : & KESIMPULAN DAN SARAN\\
\nopagebreak 
& & Bab ini berisi tentang kesimpulan dan saran dari  aplikasi yang dirancang.

\end{longtable}
%==========================================================================================================
% MULAI BAB II
%==========================================================================================================
\chapter{DASAR TEORI}
\label{chap:dasar_teori}
%==========================================================================================================
% Subbab
%==========================================================================================================

\section{\textit{Augmented Reality}}
\label{sec:AR}
\textit{Augmented reality} (AR) adalah sebuah istilah untuk lingkungan yang menggabungkan dunia nyata dan dunia virtual yang dibuat oleh komputer sehingga batas antara keduanya menjadi sangat tipis. Ronald Azuma pada tahun 1997 \cite{Azuma1997} mendefinisikan \textit{Augmented Reality} sebagai sistem yang memiliki karakteristik sebagai berikut:
\begin{itemize}
	\item Menggabungkan lingkungan nyata dan virtual
	\item Berjalan secara interaktif dalam waktu nyata
	\item Integrasi dalam tiga dimensi (3D)
\end{itemize}
Secara sederhana AR bisa didefinisikan sebagai lingkungan nyata yang ditambahkan objek virtual. Penggabungan objek nyata dan virtual dimungkinkan dengan teknologi \textit{display} yang sesuai, interaktivitas dimungkinkan melalui perangkat-perangkat input tertentu.  

AR merupakan variasi dari \textit{Virtual Environments} (VE), atau yang lebih dikenal dengan istilah \textit{Virtual Reality} (VR). Teknologi VE membuat pengguna tergabung dalam sebuah lingkungan virtual secara keseluruhan. Ketika tergabung dalam lingkungan tersebut, pengguna tidak bisa melihat lingkungan nyata di sekitarnya. Sebaliknya, AR memungkinkan pengguna untuk melihat lingkungan nyata, dengan objek virtual yang ditambahkan atau tergabung dengan lingkungan nyata. Tidak seperti VR yang sepenuhnya menggantikan lingkungan nyata, AR sekedar menambahkan atau melengkapi lingkungan nyata \cite{Azuma1997}. 

Tujuan utama dari AR adalah untuk menciptakan lingkungan baru dengan menggabungkan interaktivitas lingkungan nyata dan virtual sehingga pengguna merasa bahwa lingkungan yang diciptakan adalah nyata. Dengan kata lain, pengguna merasa tidak ada perbedaan yang dirasakan antara AR dengan apa yang mereka lihat/rasakan di lingkungan nyata. Dengan bantuan teknologi AR (seperti visi komputasi dan pengenalan objek) lingkungan nyata disekitar kita akan dapat berinteraksi dalam bentuk digital (virtual). Informasi tentang objek dan lingkungan disekitar kita dapat ditambahkan kedalam sistem AR yang kemudian informasi tersebut ditampilkan diatas \textit{layer} dunia nyata secara \textit{real-time} seolah-olah informasi tersebut adalah nyata. Informasi yang ditampilkan oleh objek virtual membantu pengguna melaksanakan kegiatan-kegiatan dalam dunia nyata. AR banyak digunakan dalam bidang-bidang seperti kesehatan, militer, industri manufaktur dan juga telah diaplikasikan dalam perangkat-perangkat yang digunakan orang banyak, seperti pada telepon genggam. 

\section{\textit{Mixed Reality}}
\label{sec:mixed_reality}
Paul Milgram and Fumio Kishino merumuskan  kerangka kemungkinan penggabungan dan peleburan dunia nyata dan dunia maya yang disebut \textit{\textbf{Milgram's Reality-Virtuality Continuum}} pada tahun  1994.  Dalam Gambar \ref{fig:milgram_continuum}, sisi yang paling kiri adalah lingkungan nyata yang hanya berisi benda nyata, dan sisi paling kanan adalah lingkungan maya yang berisi benda maya \cite{Milgram1994a}.  

\begin{figure}[h]
\begin{center}
\includegraphics{./images/Milgram_Continuum}
\caption{\label{fig:milgram_continuum} Mixed Reality}
\end{center}
\end{figure}

Dalam \textit{augmented reality}, yang lebih dekat ke sisi kiri, lingkungan bersifat nyata dan benda bersifat maya, sementara dalam \textit{augmented virtuality}, yang lebih dekat ke sisi kanan, lingkungan bersifat maya dan benda bersifat nyata. 

\section{\textit{Komponen \textit{Augmented Reality}}}
\label{sec:komponen_ar}

\section{Teknik Display \textit{Augmented Reality}}
\label{sec:AR_display}
Sistem \textit{display} AR merupakan sistem manipulasi citra yang menggunakan seperangkat optik, elektronik, dan komponen mekanik untuk membentuk citra dalam jalur optik antara mata pengamat dan objek fisik yang akan digabungkan dengan teknik AR. Bergantung kepada optik yang digunakan, citra bisa dibentuk pada sebuah benda datar atau suatu bentuk permukaan yang kompleks (tidak datar)\cite{Bimber2005}. Gambar \ref{fig:diagram_display_AR} mengilustrasikan kemungkinan citra akan dibentuk untuk mendukung AR, peletakan display bergantung dari pandangan pengguna dan objek, dan tipe citra seperti apa yang akan dihasilkan (\textit{planar} atau \textit{curved}).

\begin{figure}[h]
	\begin{center}
	\includegraphics[width=10cm]{images/diagram_display_AR}
	\caption{\label{fig:diagram_display_AR}Pembentukan citra untuk display \textit{augmented reality}}
	\end{center}
\end{figure}

Secara garis besarnya ada tiga teknik display AR \cite{Bimber2005}, yaitu sebagai berikut:

\begin{enumerate}
\item \textit{Head-Attached Display}
\item \textit{Handheld Display}
\item \textit{Spatial Display}
\end{enumerate}

\subsection {\textit{Head-Attached Display}}
\label{subsec:HAD}
\textit{Head-Attached Display} merupakan teknik display yang mengharuskan penggunanya untuk memakai sistem ini di kepala pengguna. Berdasarkan teknik citra yang terbentuk, \textit{Head-Attached Display} terbagi tiga, yaitu sebagai berikut: 
\begin{itemize}
\item \textit{Head-Mounted Display}. 
\item \textit{Head-Mounted Projectors}.
\item \textit{Virtual Retina Display}. 
\end{itemize}
Kelebihan teknik display \textit{Head-Attached Display} ini adalah lebih nyaman ke pengguna, karena citra yang terbentuk mengikuti sudut pandang pengguna.

\subsubsection {\textit{Head-Mounted Display}}
\label{subsubsec:HMD}
\textit{Head-Mounted Display} (HMD) menggabungkan citra dari objek virtual dan objek nyata dan menampilkannya langsung ke mata pengguna melalui suatu alat yang dipasang di kepala pengguna. Terdapat dua tipe utama perangkat HMD yang digunakan dalam aplikasi realitas tertambah, yaitu \textit{optical-see-through} HMD dan \textit{video see-through} HMD. Keduanya digunakan untuk berbagai jenis pekerjaan dan memiliki keuntungan dan kerugian masing-masing. Dengan \textit{optical-see-through} HMD, lingkungan nyata dilihat melalui cermin semi transparan yang diletakkan di depan mata pengguna. Cermin tersebut juga digunakan untuk merefleksikan citra yang dibentuk oleh komputer ke mata pengguna, menggabungkan lingkungan nyata dan virtual. Dengan \textit{video see-through} HMD, lingkungan nyata direkam mengunakan dua kamera video yang terintegrasi ke alat, seperti gambar \ref{fig:actual_see_trough_HMD}, dan citra yang dibentuk komputer digabung dengan video tadi untuk merepresentasikan lingkungan yang akan dilihat pengguna \cite{Rolland1994}.

\begin{itemize}
\item {\textit{Video-see-through Head-Mounted Display}}

\begin{figure}[h]
\begin{center}
\includegraphics[width=11cm]{./images/opaque_HMD}
\caption{\label{fig:opaque_HMD} Diagram \textit{Opaque} HMD}
\end{center}
\end{figure}

\textit{Video see-through} HMD bekerja dengan menggabungkan sebuah \textit{closed-view} HMD dengan satu atau dua \textsl{head-mounted} kamera video, melalui kamera video tersebut pengguna melihat ke lingkungan nyata. Video dari kamera dikombinasikan dengan citra yang dibuat oleh \textit{scene generator}, dunia nyata dan virtual digabungkan. Hasilnya dikirimkan ke monitor yang terletak di depan mata pengguna. Gambar \ref{fig:opaque_HMD} menunjukkan konsep dari \textit{Video see-through} HMD, gambar \ref{fig:actual_opaque_HMD} adalah contoh \textit{Video see-through} HMD, dengan dua video terintegrasi di bagian atas Helm \cite{Azuma1997}.

\begin{figure}[h]
\begin{center}
\includegraphics[width=5cm]{./images/actual_opaque_HMD}
\caption {\label{fig:actual_opaque_HMD} Contoh \textit{Opaque} HMD}
\end{center}
\end{figure}

\item {\textit{Optical see-through Head-Mounted Display}}

Tidak seperti penggunaan \textit{video see-through} HMD, \textit{optical see-through} HMD  menyerap cahaya dari lingkungan luar, sehingga memungkinkan pengguna untuk secara langsung mengamati dunia nyata dengan mata (gambar \ref{fig:see-trough_HMD}). Selain itu, sebuah sistem cermin yang diletakkan di depan mata pengguna memantulkan cahaya dari pencitraan grafis yang dihasilkan komputer. Pencitraan yang dihasilkan merupakan gabungan optis dari pandangan atas dunia nyata dengan pencitraan grafis \cite{Azuma1997}.

\begin{figure}[h]
\begin{center}
\includegraphics[width=11cm]{./images/see-trough_HMD}
\caption{\label{fig:see-trough_HMD} Diagram \textit{see-trough} HMD}
\end{center}
\end{figure}

\begin{figure}[h]
\begin{center}
\includegraphics[width=10cm]{./images/actual_see-trough_HMD}
\caption {\label{fig:actual_see_trough_HMD} Contoh \textit{see-through} HMD, dibuat oleh Hughes Electronics}
\end{center}
\end{figure}

\end{itemize}

\subsubsection {\textit{Head-Mounted Projectors}}
\label{subsubsec:HMDP}
\textit{Head-Mounted Projectors} Menggunakan proyektor atau panel LCD kecil dan mempunyai cahaya sendiri untuk menampilkan citra langsung ke lingkungan nyata \cite{Bimber2005}. Seperti yang ditunjukkan oleh gambar \ref{fig:hmpd}.

\begin{figure}[h]
\begin{center}
\includegraphics{./images/hmpd}
\caption {\label{fig:hmpd} Ilustrasi penggunaan dua jenis perangkat HMD yang digunakan untuk menampilkan data dan informasi tambahan}
\end{center}
\end{figure}

\subsubsection {\textit{Virtual Retina Display}}
\label{subsubsec:VRD}
\textit{Virtual retina display} (VRD), atau disebut juga dengan \textit{retinal scanning display} (RSD), memproyeksikan cahaya langsung kepada retina mata pengguna\cite{Haller2010}. VRD dapat menampilkan proyeksi citra yang penuh dan juga tembus pandang tergantung pada intensitas cahaya yang dikeluarkan, sehingga pengguna dapat menggabungkan realitas nyata dengan citra yang diproyeksikan  melalui sistem penglihatannya. VRD dapat menampilkan jarak pandang yang lebih luas daripada HMD dengan citra beresolusi tinggi\cite{Jacko2010}.  Keuntungan lain VRD adalah konstruksinya yang kecil dan ringan. Namun, VRD yang ada kini masih merupakan prototipe  yang masih terdapat dalam tahap perkembangan, sehingga masih belum dapat menggantikan HMD yang masih dominan digunakan dalam bidang AR. Gambaran sederhana VRD ini dapat dilihat pada gambar \ref{fig:VRD_diagram}.

\begin{figure}[h!]
	\centering
		\includegraphics[width=6cm]{images/virtual_retina_diagram}
	\caption{\label{fig:VRD_diagram} Diagram sederhana \textit{virtual retina display}}
\end{figure}

%\begin{figure}[h]
%\begin{center}
%\includegraphics{./images/Vrd_blocks.gif}
%\caption{\label{fig:vrd_block} Diagram Blok \textit{Virtual Retina Display}}
%\end{center}
%\end{figure}

\subsection {\textit{Handheld Display}}
\label{subsec:handheld_display}
Teknik ini menggunakan alat dengan display yang dengan mudah dapat digenggam pengguna (Tablet PC, PDA dan telepon genggam) seperti yang ditunjukkan pada gambar \ref{fig:AR_phone}. Sensor dapat berupa GPS, kompas digital ataupun kamera yang ada pada \textit{handheld} tersebut. Semua penerapan AR pada perangkat genggam menggunakan kamera untuk menggabungkan citra digital dengan lingkungan nyata, \textit{Handheld} AR sangat menjanjikan untuk tujuan komersial. Dua kelebihan utama dari \textit{Handheld} AR adalah mobilitas perangkat yang mudah dan salah satu perangkat genggam yang banyak digunakan (telepon genggam) telah banyak dilengkapi kamera.

\begin{figure}[h!]
	\centering
		\includegraphics[width=8cm]{images/handphone_AR}
	\caption{\label{fig:AR_phone} Contoh augmented reality dengan \textit{handphone}}
\end{figure}

\subsection {\textit{Spatial Display}}
\label{subsec:spatial_display}
Dalam \textit{Spatial Augmented Reality} (SAR), objek nyata digabungkan langsung dengan citra yang terintegrasi langsung ke lingkungan nyata. Contohnya, citra diproyeksikan ke lingkungan nyata menggunakan proyektor digital atau tergabung dengan lingkungan menggunakan panel display \cite{Ramesh1998}. Perbedaan utama pada SAR dibanding teknik display sebelumnya adalah displaynya terpisah dengan pengguna. SAR memiliki kelebihan dari HMD dan handheld, sistem ini bisa digunakan oleh banyak orang pada waktu bersamaan tanpa perlu mengenakan suatu alat.

Ada tiga teknik display dalam SAR \cite{Bimber2005}, yaitu sebagai berikut:
\begin{enumerate}
\item \textit{Screen-Based Video See-Through Displays}\\
\textit{Screen-based }AR menggabungkan citra dan lingkungan nyata yang ditampilkan ke sebuah monitor, seperti yang ditunjukkan pada gambar \ref{fig:SAR_screen}.

\begin{figure}[h]
	\centering
		\includegraphics[width=6cm]{images/screen_SAR}
	\caption{\label{fig:SAR_screen} Contoh \textit{Screen-Based Video See-Through Displays} }
\end{figure}

\item \textit{Spatial Optical See-Through Displays}\\
Sistem ini menghasilkan citra yang ditampilkan langsung ke lingkungan nyata. Komponen yang penting dalam sistem ini meliputi \textit{spatial optical combiners} (\textit{planar} atau \textit{curved beam combiners}), layar transparan atau hologram.

\item \textit{Projection-Based Spatial Displays}\\
Sistem ini memproyeksikan citra secara langsung pada permukaan objek fisik daripada menampilkannya pada sebuah bidang pencitraan dalam penglihatan pengguna. Sistem ini menggunakan banyak proyektor yang digunakan untuk meningkatkan wilayah tampilan serta meningkatkan kualitas citra.

\end{enumerate}

\section{Model Tiga Dimensi (3D)}
\label{sec:model_3d}
Pemodelan Tiga Dimensi (3D) (3D \textit{modeling} atau dikenal juga dengan \textit{meshing}) adalah proses pembuatan representasi matematis permukaan tiga dimensi dari suatu objek dengan \textit{software} tertentu. Produk hasil pemodelan itu disebut model 3D. Model 3D tersebut dapat ditampilkan sebagai citra dua dimensi melalui sebuah proses yang disebut \textit{3D rendering}.
%
%Models may be created automatically or manually. The manual modeling process of preparing geometric data for 3D computer graphics is similar to plastic arts such as sculpting.
%
Model 3D direpresentasikan dari kumpulan titik dalam 3D, terhubung oleh berbagai macam entitas geometri, seperti segitiga, garis, permukaan lengkung, dan lain sebagainya. Berdasarkan hal tersebut, model 3D bisa dibuat manual (seperti seni memahat), secara algoritma (pemodelan prosedural), atau \textit{scanning}. Hasil akhir dari citra 3D adalah sekumpulan poligon. Model dengan jumlah poligon yang lebih banyak memerlukan waktu yang lebih lama untuk di-\textit{render} oleh komputer, karena setiap permukaan memiliki tekstur dan \textit{shading} tersendiri. 

\section{Aplikasi Komputer Berbasis Web (\textit{Rich Internet Application})}
\label{sec:aplikasi_web}

Pada tahun 90-an, "mengakses \textit{web}" berarti mengakses tulisan dan gambar statis secara \textit{online} . Sejalan dengan berkembangnya koneksi internet, kebutuhan akan suatu konten yang lebih kaya, responsif juga meningkat. Pada tahun 2002, Macromedia menemukan istilah \textit{Rich Internet applications} (RIAs). RIAs menggabungkan fleksibilitas, tingkat responsif, dan kemudahan aplikasi berbasis desktop. RIAs atau dikenal dengan aplikasi berbasis \textit{web} adalah aplikasi yang mempunyai karakteristik seperti aplikasi \textit{desktop}, biasanya didistribusikan atau diakses lewat \textit{browser web} standar, menggunakan \textit{web plug-in}. Contoh RIAs meliputi Ajax, Curl, GWT, Adobe Flash/Adobe Flex/AIR, Java/JavaFX, Mozilla's XUL dan Microsoft Silverlight. 

Aplikasi berbasis \textit{web} semakin banyak dikembangkan, karena kelebihan aplikasi \textit{web} yang sangat baik untuk aplikasi dengan model \textit{client-server}. Beberapa kelebihan aplikasi \textit{web} dibanding aplikasi \textit{desktop}:
\begin{enumerate}
	\item \textit{Running anywhere} (berjalan/beroperasi dimana saja), cukup instalasi pada satu \textit{server} dan tanpa perlu instalasi apapun selama ada \textit{web browser} di sisi \textit{client}, pengguna langsung  dapat menjalankan aplikasi tersebut tanpa konfigurasi apapun. Secara default, \textit{web browser} pada sistem operasi apapun telah tersedia sehingga kebutuhan akan \textit{web browser} untuk menjalankan aplikasi bukanlah suatu kendala berarti.
	\item \textit{Easy to update} (mudah untuk diubah/diperbaharui), cukup \textit{update} pada sisi \textit{server} maka aplikasi pada \textit{client} akan langsung menggunakan versi ter-\textit{update} tanpa harus instalasi dulu pada \textit{client}.
	\item \textit{Requirement} (kebutuhan)pada \textit{client} tidak terlalu besar. Karena \textit{running} aplikasi bersifat \textit{stateless} dan pada sisi \textit{client} hanya sebagai \textit{interface} maka spesifikasi hardware pada \textit{client} tidak harus canggih. 
	\item Tampilan yang lebih beragam.
\end{enumerate} 

\section {Adobe Flash \textit{Platform}}
\label{subsec:adobe_flash_platform}
Kebanyakan \textit{designer} dan \textit{developer} menggunakan Adobe Flash ataupun Adobe Flex, yang merupakan bagian dari platform Adobe Flash, untuk mengembangkan RIAs. Flash merupakan suatu \textit{environment} untuk membuat konten yang interaktif dan kaya fitur dalam dunia \textit{web}. Begitu juga Flex merupakan sebuah \textit{framework} \textit{cross-platform} untuk mengembangkan RIAs. Konten yang dibuat dengan Flash dan Flex di-\textit{deploy} menggunakan Adobe Flash Player

\subsection {Adobe Flash}
\label{subsec:adobe_flash}
Adobe Flash (dulunya Macromedia Flash) adalah platform multimedia yang aslinya dibuat oleh Macromedia dan saat ini dikembangkan dan didistribusikan oleh Adobe Systems. Sejak pengenalannya pada Tahun 1996, Flash telah menjadi metode yang popular untuk menambahkan animasi dan interaktivitas ke halaman web. Komponen Flash untuk mengintegrasikan video ke halaman web, dan yang terbaru saat ini, untuk mengembangkan RIAs.

Flash dapat memanipulasi vector dan raster grafik, serta mendukung \textit{streaming} dua arah audio dan video. Flash menggunakan bahasa \textit{script} yang disebut Action Script. Banyak produk \textit{software}, sistem dan \textit{device} dapat menampilkan konten flash, contohnya Adobe Flash player, yang tersedia gratis bagi sebagian besar \textit{web browser}. Beberapa ponsel dan alat elektronik lainnya juga dapat menampilkan konten Flash, menggunakan Flash lite.

\textit{File} dalam format SWF, biasanya disebut "ShockWave Flash movies", "Flash movies" atau "Flash games", yang biasanya memiliki sebuah ekstensi .swf dan dapat menjadi objek di halaman \textit{web}. \textit{File} tersebut pada dasarnya dijalankan dengan Flash Player itu sendiri atau digabungkan dengan "Projector" (\textit{video flash} yang dapat berjalan sendiri dengan ekstensi .exe di Microsoft Windows atau .hqx untuk Macintosh). \textit{File} Flash Video memiliki ekstensi .flv dan juga digunakan dalam .swf atau dijalankan melalui aplikasi yang dapat menjalankan \textit{file} flv.

\subsection {Adobe Flex}
\label{subsec:adobe_flex}
Adobe Flex adalah paket pengembangan \textit{software} yang dirilis oleh Adobe systems untuk pengembangan dan aplikasi \textit{cross-platform rich internet} berbasis platform adobe flash. Aplikasi flex dapat ditulis menggunakan Adobe Flash Builder sebagai IDE (\textit{Integrated Development Evironment}) dan menggunakan Flex \textit{Software Development Kit} (SDK) yang tersedia gratis dari Adobe. 

Flex membuat \textit{workflow} dan model pemrograman yang familiar terhadap \textit{developer} aplikasi berbasis \textit{web}. FLex menggunakan MXML dan ActionScript. MXML merupakan sebuah bahasa yang berbasis XML, menawarkan cara membangun dan menata \textit{user interface}. Interaktivitas dicapai melalui penggunaan ActionScript, yaitu bahasa utama Flash Player.

Flex menyediakan dua kompiler: mxmlc dan compc. Kompiler compc dan mxmlc bisa dijalankan dari Flex Builder atau pun dari \textit{command line}. compc digunakan untuk mengkompilasi komponen, kelas, and file lainnya ke file SWC atau RSL yang digunakan sebagai \textit{library} untuk aplikasi. mxmlc untuk mengkompilasi ActionScript dan SWC ke \textit{file} SWF, seperti ditunjukkan pada gambar \ref{fig:compile_flex}. Setelah aplikasi dikompilasi dan di-\textit{deploy} ke \textit{web server}, pengguna bisa mengakses file SWF tersebut dan dijalankan via \textit{web browser}.

\begin{figure}
\begin{center}
\includegraphics[width=9cm]{./images/compile_flex.png}
\caption{\label{fig:compile_flex} Diagram Kompilasi Flex}
\end{center}
\end{figure}

\subsection {ActionScript}
\label{subsec:action_script}
ActionScript merupakan bahasa pemrograman berorientasi objek yang berdasarkan ECMAScript (bahasa yang distandarisasi oleh Ecma International dalam spesifikasi ECMA-262 dan ISO/IEC 16262). ActionScript  terutama digunakan untuk pengembangan \textit{website} dan \textit{software} menggunakan Adobe Flash Player (dalam bentuk \textit{file} SWF yang diintegrasikan  ke halaman \textit{web}), ActionScript juga digunakan pada beberapa aplikasi untuk \textit{database}  (seperti Alpha Five). ActionScript pada awalnya didesain untuk mengatur animasi vektor 2D sederhana yang dibuat di Adobe Flash, dengan berkembangnya. Versi terakhir dari ActionScript menambahkan kemungkinan penggunaan  untuk pembuatan \textit{web} berbasis \textit{game} dan RIAs dengan media \textit{streaming} (seperti video dan audio).

\section{\textit{Library} Pendukung Augmented Reality Pada Platform Flash}
\label {sec:library_AR}
\textit{Library} atau Pustaka, dalam ilmu komputer adalah koleksi dari rutin-rutin program yang digunakan untuk membangun dan mengembangkan \textit{software}. \textit{Library}, umumnya mengandung kode program dan data pembantu (banyak \textit{programmer} menyebutnya sebagai \textit{helper}), yang menyediakan layanan-layanan kepada program-program independen. Penggunaan \textit{library} sangat diperlukan untuk mengembangkan aplikasi AR secara cepat. \textit{Library} AR yang dikenal luas dalam platform Flash adalah FLARToolKit.

\subsection {FLARToolKit}
\label {subsec:FLARToolKit}
%\subsection {FLARToolKit}
%\label{subsec:FLARToolKit}
%Pada bulan november 2008, sebuah grup pemrograman di Jepang, mengembangkan banyak proyek dengan ActionScript sehingga memberi ide kepada para \textit{developer} akan apa yang bisa dilakukan bahasa tersebut. FLARToolKit, pertama kali dikembangkan oleh Tomohiko Koyama (aka Saqoosha), ia mengenalkan AR ke dunia \textit{web}, dan ke segmen pengguna yang lebih luas.

FLARToolKit adalah \textit{tracking system library} yang bersifat \textit{open-source} sehingga memungkinkan \textit{programmer} dengan mudah mengembangkan aplikasi AR, FLARToolKit merupakan \textit{porting} (perubahan terhadap \textit{software} untuk menjadikannya dapat digunakan di lingkungan yang berbeda) yang paling terakhir dari ARToolkit, yaitu sebuah libary AR C++ yang awalnya dikembangkan oleh Dr. Hirokazu Kato di  Human Interface Technology Lab University of Washington. Dengan datangnya ActionScript 3.0, para pengembang seperti Mario Klingemann dan lainnya mulai bereksperimen dengan teknik analisis image secara \textit{real-time} untuk  Flash Player. Saqoosha meneruskan hal ini, dan mem-\textit{porting} FLARToolKit dari NYARToolkit (sebuah Java/C-sharp/Android \textit{port} dari ARToolkit).

FLARToolKit hanya merupakan \textit{library} untuk \textit{tracking} pada AR, untuk menampilkan objek 3D di lingkungan Flash, FLARToolKit memerlukan sebuah \textit{library} 3D. Beberapa \textit{library} 3D yang didukung oleh FLARToolKit adalah sebagai berikut:
\begin{itemize}
\item Alternativa3D
\item Away3D
\item Away3D Lite
\item Papervision3D
\item Sandy3D
\end{itemize}

Proses FLARToolKit secara garis besarnya sebagai berikut: 

\begin{enumerate}
	\item Mengambil video dari \textit{webcam}.
	\item Binarisasi citra masukan(\textit{thresholding}).
	\item Memberi penanda (\textit{labelling}).
	\item Deteksi area persegi (\textit{Marker Outline Detection}). 
	\item Pencocokan pola.
	\item Menghitung transform matrix.
	\item Merender obyek 3D.
\end{enumerate}

\subsubsection {Mengambil Video dari \textit{Webcam}}
\label{subsubsec:capture_webcam}
Langkah awal yang harus dilakukan adalah mendapatkan masukan video dari sebuah \textit{webcam}, seperti yang ditunjukkan gambar \ref{fig:capture_from_webcam}. Video yang di-\textit{streaming} secara \textit{real-time} ini akan diolah oleh sistem untuk dianalisa \textit{frame} per \textit{frame}. Sebelum \textit{webcam} digunakan, \textit{webcam} harus dikaliberasi terlebih dahulu. Kaliberasi \textit{webcam} merupakan bagian yang sangat penting dalam proses pengambilan masukan video. Hal ini disebabkan oleh distorsi pada lensa \textit{webcam} yang tiap-tiap kamera berbeda karakteristiknya (gambar \ref{fig:distorted_image}). Tujuan dari kalibrasi \textit{webcam} adalah untuk menghitung tingkat distorsi dari sebuah lensa \textit{webcam} yang digunakan agar citra yang dihasilkan mendekati citra ideal.

\begin{figure}[h]
\begin{center}
\includegraphics[width=7cm]{./images/flartk/capture_webcam.JPG}
\caption{\label{fig:capture_from_webcam} Mengambil citra dari \textit{webcam}}
\end{center}
\end{figure}

%Menggunakan \textit{class} kamera dan video kemudian menggambar video contohnya ke BitmapData. 

\begin{figure}[h]
\begin{center}
\includegraphics[width=6cm]{./images/flartk/distorted_image.JPG}
\caption{\label{fig:distorted_image} Perbandingan antara citra yang ideal dengan citra yang disebabkan oleh faktor distorsi}
\end{center}
\end{figure}

\subsubsection {Binarisasi Citra Masukan (\textit{Thresholding})}
\label{subsubsec:thresholding}
Langkah pertama pada aplikasi visi komputer yang terletak pada deteksi tepi adalah untuk men-\textit{threshold} sumber citra atau disebut juga binarisasi seperti yang ditunjukkan gambar \ref{fig:thresholding}. \textit{Thresholding} mengkonversi citra ke citra binari sehingga memudahkan untuk komputasi. Sebuah citra binari dibuat dengan mengubah pixel yang lebih cerah daripada nilai \textit{threshold} ke suatu warna, dan pixel yang lebih gelap daripada nilai \textit{threshold} ke suatu warna lainnya (didefinisikan sebagai \textit{gray-scale} atau hitam-putih). Nilai \textit{threshold} berada pada angka 0 – 255 dan secara default, \textit{threshold} bernilai 100. Fungsi dari proses ini adalah untuk membantu sistem agar dapat mengenali bentuk segi empat dan pola di marker pada citra yang diterima. Nilai \textit{threshold} dapat dirubah dan disesuaikan dengan kondisi cahaya disekitar \textit{marker} untuk tetap membuat \textit{marker} terlihat sebagai segi empat, karena ketika cahaya disekitar \textit{marker} berkurang ataupun berlebih pada saat proses \textit{thresholding}, sistem tidak dapat mendeteksi \textit{marker}.
%\begin{figure}[h]
%\begin{center}
%\includegraphics[width=7cm]{./images/thresholding}
%\caption{\label{fig:thresholding} Thresholding}
%\end{center}
%\end{figure}
\begin{figure}
\begin{center}
\includegraphics[width=7cm]{./images/flartk/binarize.JPG}
\caption{\label{fig:thresholding} Thresholding}
\end{center}
\end{figure}

\subsubsection {Memberi Penanda (\textit{Labeling})}
\label{subsubsec:labeling}
Langkah berikutnya dari FLARToolKit adalah menemukan area yang berdampingan dalam citra yang di-\textit{treshold}, khususnya dalam area dibawah \textit{threshold} (area yang lebih gelap). Area yang berdampingan diberi tanda dengan warna yang berbeda dengan tujuan untuk mengidentifikasi area, proses \textit{labelling} dapat dilihat pada gambar \ref{fig:labelling}.
%Dengan menggunakan fungsi \textit{BitmapData.getColorBoundsRect} dan \textit{BitmapData.foodFill}, 

\begin{figure}[h]
\begin{center}
\includegraphics[width=7cm]{./images/flartk/labelling.JPG}
\caption{\label{fig:labelling} Setiap area putih ditandai dengan warna yang berbeda.}
\end{center}
\end{figure}

\subsubsection {Deteksi Area Persegi (\textit{Marker Outline Detection})}
\label{subsubsec:outline_detection}
Langkah selanjutnya, FLARToolKit mencari area yang kemudian ditandai sebagai persegi \textit{(marker outline)}. Setelah citra mengalami proses \textit{thresholding} dan \textit{labelling}, FLARToolkit akan mengenali bentuk dan pola yang ada pada marker. FLARToolkit akan mencari bagian yang memiliki bentuk segi empat dan menandainya. FLARToolKit juga akan menghilangkan area yang tidak berbentuk segi empat sehingga yang akan ditampilkan pada layar hanyalah area yang memiliki bentuk segi empat (gambar \ref{fig:find_squares}.
%With candidates for marker locations, FLARToolKit then proceeds to search the labeled areas for shapes that could be transformed squares (i.e. marker outlines).

\begin{figure}[h]
\begin{center}
\includegraphics[width=7cm]{./images/flartk/find_squares.JPG}
\caption{\label{fig:find_squares} Mencari area persegi (Marker Outline Detection)}
\end{center}
\end{figure} 

\subsubsection {Pencocokan Pola}
\label{subsubsec:pattern_matching}
Setelah semua area persegi ditandai, FLARToolKit menganalisa citra yang berada di dalam persegi dan membandingkan polanya dengan sekumpulan pola yang telah ditentukan (pencocokan pola). FLARToolKit mengekstrak pola di dalam persegi menggunakan transformasi \textit{homography}. FLARToolKit memberikan sebuah nilai '\textit{confidence}' kepada setiap pola yang cocok, jika kecocokannya di atas nilai yang telah ditentukan maka polanya dinyatakan cocok. 

\begin{figure}
\begin{center}
\includegraphics[width=7cm]{./images/flartk/marker_spec.JPG}
\caption{\label{fig:marker_spec} Spesifikasi pola marker}
\end{center}
\end{figure} 

Spesifikasi pola marker (gambar \ref{fig:marker_spec}):
\begin{itemize}
	\item Harus berupa persegi.
	\item Hanya 50\% dari tengah area yang digunakan untuk proses pencocokan pola.
	\item Pola marker secara \textit{default}-nya adalah 16 x 16 titik.
	\item Ukuran pola bisa lebih besar, tapi membutuhkan waktu yang lebih lama untuk diproses.
\end{itemize}

%\begin{figure}[h!]
%	\begin{center}
%		\subfloat {\includegraphics[width=5cm]{./images/flartk/matching_pattern.JPG} \label{fig:pattern_matching1}}\\
%		\subfloat {\includegraphics[width=5cm]{./images/flartk/matching_pattern2.JPG} \label{fig:pattern_matching2}}\\
%		\subfloat {\includegraphics[width=5cm]{./images/flartk/matching_pattern3.JPG} \label{fig:pattern_matching3}}\\
%		\subfloat {\includegraphics[width=5cm]{./images/flartk/matching_pattern4.JPG} \label{fig:pattern_matching4}}
%		\caption {\label{fig:pattern_matching} Proses pencocokan pola}
%	\end{center}
%\end{figure}

\subsubsection {Menghitung Transformasi Matriks}
\label{subsubsec:calculate_transform_matrix}
Transformasi matriks dihitung dari titik-titik persegi marker yang dideteksi. Matriks tersebut digunakan untuk proses \textit{render} objek 3D.% Algoritmanya dapat dilihat dalam paper di http://www.hitl.washington.edu/artoolkit/publications/.
%
%\begin{figure}[h!]
%\begin{center}
%\includegraphics[width=7cm]{./images/flartk/calculate_transform_matrix.JPG}
%\caption{\label{fig:calculate_matrix} Kalkulasi transformasi matriks}
%\end{center}
%\end{figure} 

\subsubsection {Me-\textit{render} Objek 3D}
\label{subsubsec:render_the_3d_objects}
FLARToolKit menggunakan transformasi matriks yang dikalkulasikan di step sebelumnya dan menampilkan objek yang sesuai dengan sebuah \textit{library} 3D, seperti yang ditunjukkan gambar \ref{fig:render_3d_objects}. FLARToolKit menyertakan kelas pendukung yang mengkonversikan transformasi matriks FLARTollKit ke setiap kelas matriks internal \textit{library} 3D tersebut.

\begin{figure}[h!]
\begin{center}
\includegraphics[width=7cm]{./images/flartk/render_3d_objects.JPG}
\caption{\label{fig:render_3d_objects} Render objek 3D}
\end{center}
\end{figure} 

\subsection {Papervision3D}
\label{subsec:Papervision3D}
Salah satu \textit{library} 3D (biasa disebut juga dengan \textit{engine} 3D) untuk platform Flash adalah Papervision3D, Papervison3D merupakan \textit{library} 3D yang bersifat \textit{open source}. Dengan Papervision 3D, model 3D bisa dibuat di platform Flash menggunakan kelas-kelas dalam ActionScript ataupun juga dengan meng-\textit{import} model yang dibuat dari \textit{software} pemodelan 3D. Model yang telah dibuat atau di-\textit{import}, bisa dikendalikan dengan ActionScript.

\subsection {FLARManager}
\label{subsec:FLARManager}
Visi komputer dalam konteks \textit{web} memiliki banyak kesulitan untuk diterapkan. Masalah utama ialah kurangnya kendali terhadap kondisi lingkungan \textit{end-user}. Kurangnya atau ketidakteraturan pencahayaan membuat analisis dari sebuah citra sangat sulit, dan permasalahan tersebut berdampak pada FLARToolKit. Selain itu, FLARToolKit sangat sulit digunakan karena minimnya dokumentasi, dan \textit{developer} harus membuat sebuah \textit{manager} untuk mengatur konfigurasi kamera, marker dan juga objek virtual yang akan ditampilkan. 

FLARManager adalah sebuah \textit{framework} sederhana yang memudahkan untuk membuat aplikasi AR dengan Flash terutama yang menggunakan \textit{libray} FLARToolKit. \textit{Framework} merupakan kumpulan komponen kelas sekaligus kerangka dalam pemrograman yang memudahkan \textit{programmer} untuk membuat suatu aplikasi yang siap pakai. FLARManager bisa meningkatkan akurasi dan reabilitas dari proses dengan hanya merubah konfigurasinya yang berupa file dengan format xml. FLARManager dapat mendeteksi lebih dari satu \textit{marker} dan pola dalam waktu bersamaan, sehingga bisa menampilkan objek virtual lebih dari satu.

FLARToolkit menyediakan akses ke sumber citra yang diprosesnya, dan juga hasil dari FLARToolkit setelah di analisa. Dengan demikian, FLARManager bisa difokuskan pada fungsionalitas dalam deteksi dan \textit{tracking} dari \textit{marker}. FLARManager sebisa mungkin menghindari perubahan terhadap FLARToolkit. Dengan tidak tergantungnya FLARManager terhadap FLARToolkit, setiap proyek bisa terus dikembangkan secara terpisah, dan FLARManager secara teori bisa diterapkan pada \textit{libray} Flash AR lainnya.

Sampai saat tulisan ini dibuat, FLARManager telah mendukung \textit{tracking library} sebagai berikut :
\begin{itemize}
\item FLARToolkit
\item flare*tracker
\item flare*NFT
\end{itemize}



%==========================================================================================================
% MULAI BAB III
%==========================================================================================================
\chapter{PERANCANGAN DAN IMPLEMENTASI}
\label{chap:implementasi}
%==========================================================================================================
% Subbab
%==========================================================================================================
\section{Teknologi}
\label{sec:teknologi_implementasi}

\subsection {\textit{Rich Internet Application} (RIA)}
\label{subsec:RIA}
Rich Internet applications (RIAs) are web applications that have most of the characteristics of desktop applications, typically delivered either by way of a standards based web browser, via a browser plug-in, or independently via sandboxes or virtual machines.[1] Examples of RIA frameworks include Ajax, Curl, GWT, Adobe Flash/Adobe Flex/AIR, Java/JavaFX,[2], Mozilla's XUL and Microsoft Silverlight.[3]

The term was introduced in March 2002 by vendors like Macromedia who were addressing limitations at the time in the "richness of the application interfaces, media and content, and the overall sophistication of the solutions" by introducing proprietary extensions.[4] There has been some debate about whether Ajax properly qualifies as an RIA or whether the term should be reserved for plug-in based technologies. However Ajax clearly demonstrates all the core characteristics of an RIA and current opinion appears to hold that it should be therefore be included in this category[5].

As web standards have developed and the compliance of web browser has improved, the need for plug-in based RIAs has diminished. The rapid evolution of Javascript and the emergence of a broad range of Ajax-based desktop-like widget sets have continued this trend. HTML 5 takes this even further by delivering an extensive pseudo-application platform. [6]

With a few but growing number of exceptions (most notably YouTube which currently relies on Adobe Flash for video playback) the vast majority of the most popular web sites are native web applications. However, many major sites make use of RIA frameworks such as Ajax and Adobe Flash. Online gaming is an area where plugin-based RIAs are particularly prevalent. Applications (such as Dimdim) which require access to video capture also tend to use RIAs (with the notable exception of Gmail which uses its own task-specific browser plug-in[7]).

\subsubsection {Adobe Flash}
\label{subsubsec:adobe_flash}
Adobe Flash (formerly Macromedia Flash) is a multimedia platform originally acquired by Macromedia and currently developed and distributed by Adobe Systems. Since its introduction in 1996, Flash has become a popular method for adding animation and interactivity to web pages. Flash is commonly used to create animation, advertisements, and various web page Flash components, to integrate video into web pages, and more recently, to develop rich Internet applications.

Flash can manipulate vector and raster graphics, and supports bidirectional streaming of audio and video. It contains a scripting language called ActionScript. Several software products, systems, and devices are able to create or display Flash content, including Adobe Flash Player, which is available free for most common web browsers, some mobile phones and for other electronic devices (using Flash Lite). The Adobe Flash Professional multimedia authoring
program is used to create content for the Adobe Engagement Platform, such as web applications, games and movies, and content for mobile phones and other embedded devices.

Files in the SWF format, traditionally called "ShockWave Flash" movies, "Flash movies" or "Flash games", usually have a .swf file extension and may be an object of a web page, strictly "played" in a standalone Flash Player, or incorporated into a Projector, a self-executing Flash movie (with the .exe extension in Microsoft Windows or .hqx for Macintosh). Flash Video files have a .flv file extension and are either used from within .swf files or played through a flv-aware player, such as VLC, or QuickTime and Windows Media Player with external codecs added.

\subsubsection {Adobe Flex}
\label{subsubsec:adobe_flash}
Adobe Flex is a software development kit released by Adobe Systems for the development and deployment of cross-platform rich Internet applications based on the AdobeFlash platform. Flex applications can be written using Adobe Flex Builder or by using the freely available Flex compiler from Adobe.

Traditional application programmers found it challenging to adapt to the animation metaphor upon which the Flash Platform was originally designed. Flex seeks to minimize this problem by providing a workflow and programming model that is familiar to these developers. MXML, an XML-based markup language, offers a way to build and lay out graphic user
interfaces. Interactivity is achieved through the use of ActionScript, the core language of Flash Player that is based on the ECMAScript standard.

The Flex SDK comes with a set of user interface components including buttons, list boxes, trees, data grids, several text controls, and various layout containers. Charts and graphs are available as an add�]on. Other features like web services, drag and drop, modal dialogs, animation effects, application states, form validation, and other interactions round out the application framework.

In a multitiered model, Flex applications serve as the presentation tier. Unlike pagebased HTML applications, Flex applications provide a stateful client where significant changes to the view don't require loading a new page. Similarly, Flex and Flash Player provide many useful ways to send and load data to and from server-side components without requiring the client to reload the view. Though this functionality offered advantages over HTML and JavaScript development in the past, the increased support for XMLHttpRequest in major browsers has made asynchronous data loading a common practice in HTML-based development as well.

\subsubsection {Action Script}
\label{subsubsec:action_script}
ActionScript is a scripting language based on ECMAScript. ActionScript is used primarily for the development of websites and software using the Adobe Flash Player platform (in the form of SWF files embedded into Web pages), but is also used in some database applications (such as Alpha Five), and in basic robotics, as with the Make Controller Kit. Originally developed by Macromedia, the language is now owned by Adobe (which acquired Macromedia in 2005). ActionScript was initially designed for controlling simple 2D vector animations made in Adobe Flash (formerly Macromedia Flash). Later versions added functionality allowing for the creation of Web-based games and rich Internet applications with streaming media (such as video and
audio).

\subsubsection {Adobe Flash Builder}
\label{subsubsec:adobe_flash_builder}
Flash Builder 4 is Adobe's professional Flex IDE built on Eclipse. It can run either as astandalone tool or as a plugin to an existing Eclipse installation. Flex SDK includes the Flex framework (also known as the Flex class library), Flex command-line compilers, Adobe Integrated Runtime (AIR) framework, Adobe AIR command-line compilers, the Flex debugger,
the ASDoc utility, and the debugger version of Flash Player. Use the Flex SDK to develop, compile, and deploy Flex applications that connect to XML and SOAP web services with no additional charges or server licensing required. This release also includes Data Visualization features.

\subsection {ARToolKit}
\label{subsec:artoolkit}
ARToolKit is a C and C++ language software library that lets programmers easily develop Augmented Reality applications. ARToolKit was originally developed by Dr. Hirokazu Kato, and its ongoing development is being supported by the Human Interface Technology Laboratory (HIT Lab) at the University of Washington, HIT Lab NZ  at the University of Canterbury, New Zealand, and ARToolworks, Inc, Seattle.

One of the most difficult parts of developing an Augmented Reality application is precisely calculating the user's viewpoint in real time so that the virtual images are exactly aligned with real world objects. ARToolKit uses computer vision techniques to calculate the real camera position and orientation relative to marked cards, allowing the programmer to overlay virtual objects onto these cards. The fast, precise tracking provided by ARToolKit should enable the rapid development of many new and interesting AR applications. \sout{ARToolKit uses computer vision algorithms to solve this problem. The ARToolKit video tracking libraries calculate the real camera position and orientation relative to physical markers in real time. This enables the easy development of a wide range of Augmented Reality applications.} Some of the features of ARToolKit include:

\begin{itemize}
	\item Single camera position/orientation tracking.
  \item Tracking code that uses simple black squares.
  \item The ability to use any square marker patterns. 
  \item Easy camera calibration code.
  \item Fast enough for real time AR applications.
  \item SGI IRIX, Linux, MacOS and Windows OS distributions.
  \item Distributed with complete source code.
\end{itemize}

ARToolKit currently runs on the SGI IRIX, PC Linux, Mac OS X, and PC Windows (95/98/NT/2000/XP) operating systems. The last version of ARToolKit is completly multi-platform. The functionality of each version of the toolkit is the same, but the performance may vary depending on the different hardware configurations. 

\subsubsection{Prinsip Kerja ARToolKit}
\label{subsubsec:prinsip_artoolkit}
ARToolKit applications allow virtual imagery to be superimposed over live video of the real world. Although this appears magical it is not. The secret is in the black squares used as tracking markers. The ARToolKit tracking works as follows:
   
\begin{enumerate}
	\item The camera captures video of the real world and sends it to the computer.
	\item Software on the computer searches through each video frame for any square shapes
	\item If a square is found, the software uses some mathematics to calculate the position of the camera relative to the black square.
	\item Once the position of the camera is known a computer graphics model is drawn from that same position.
	\item This model is drawn on top of the video of the real world and so appears stuck on the square marker.
	\item The final output is shown back in the handheld display, so when the user looks through the display they see graphics overlaid on the real world.
\end{enumerate}

\begin{figure}
	\centering
		\includegraphics[width=14cm]{images/diagram_artoolkit.jpg}
	\caption{Diagram ARToolKit}
	\label{fig:diagram_artoolkit}
\end{figure}

The figure \ref{fig:diagram_artoolkit} summarizes these steps. ARToolKit is able to perform this camera tracking in real time, ensuring that the virtual objects always appear overlaid on the tracking markers. 

\subsection {FLARToolKit}
\label{subsec:FLARToolKit}
Back in November 2008, a group of Japanese coders, working largely under the radar, unveiled a project that redefined many ActionScript developers' ideas of what the language could do. FLARToolKit, developed primarily by Tomohiko Koyama (aka Saqoosha), introduced  augmented reality to the web, and to a large segment of the population as a whole.

FLARToolKit is the latest in a series of ports of ARToolkit \sout{,an augmented reality C++ library originally developed by Dr. Hirokazu Kato at the Human Interface Technology Lab at University of Washington}. With the advent of ActionScript 3.0, developers like Mario Klingemann and others began experimenting with realtime image analysis techniques for Flash Player. Saqoosha picked up on this, and ported FLARToolKit from NYARToolkit, a Java/C-\sout{sharp}/Android port of ARToolkit. 

FLARToolKit made its biggest initial splash at the hands of North Kingdom, the Swedish interactive agency that developed GE's SmartGrid augmented reality campaign. Since then, a host of AR applications have made their way to the web via FLARToolKit; the majority of them are variations on the theme of 3D characters dancing on top of live video, or games. As time goes on, however, creative developers will imagine new, creative, and useful applications of the technology.

\subsection {The FLARToolKit process}
\label{subsec:FLARToolKit}
FLARToolKit does the following steps to create final AR image.
\begin{enumerate}
	\item Capture the image from webcam.
	\item Binarize input image (thresholding).
	\item Labeling.
	\item Find squares.
	\item Matching with patterns.
	\item Calculate transform matrix.
	\item Render the 3D objects.
\end{enumerate}

\subsubsection {\textit{Thresholding}}
\label{subsubsec:thresholding}
The first step in many computer vision applications that rely on edge detection is to  threshold the source image. A binary image is made by changing pixels brighter than a threshold value to one color, and pixels darker than the threshold to another. 
\begin{figure}[h]
\begin{center}
\includegraphics[width=7cm]{./images/thresholding}
\caption{\label{fig:thresholding} Thresholding}
\end{center}
\end{figure}
Thresholding separates the source image into a binary image, making analysis less computationally expensive.

\subsubsection {\textit{Labeling}}
\label{subsubsec:labeling}
FLARToolKit's next step is to find contiguous areas in the thresholded image, speficially within the areas below the threshold (darker areas). Using BitmapData.getColorBoundsRect and BitmapData.floodFill, contiguous areas are 'labeled' with unique colors, used later to id the areas.
\begin{figure}[h]
\begin{center}
\includegraphics[width=7cm]{./images/labeling}
\caption{\label{fig:labeling} Each contiguous area of white (corresponding to dark areas of the source image) is 'labeled' with a different color. }
\end{center}
\end{figure}

\subsubsection {Marker Outline Detection}
\label{subsubsec:labeling}
With candidates for marker locations, FLARToolKit then proceeds to search the labeled areas for shapes that could be transformed squares (i.e. marker outlines). 

\subsubsection {Pattern Matching}
\label{subsubsec:pattern_matching}
Once all marker outline possibilities have been established, FLARToolKit analyzes the areas of the image within the outlines and compares the contents with the list of patterns the developer has asked FLARTookit to detect. FLARToolKit assigns a 'confidence' value to all of the matches; matches that are at or above the confidence level specified by the developer are reported as pattern matches. 

\subsection {Papervision 3D}
\label{subsec:papervision3D}
Papervision3D is an open source 3D engine for the Flash platform.

\section {Augmented Reality (AR) menggunakan FLARToolKit dan Papervision3D}
\label{sec:AR_flar_PV3D}
AR is combining real world objects with virtual  ones (computer generated). Major experimentation with AR has been in field of video processing, where live video feeds are combined with 2D or 3D objects. In Flash, this is usually done with a webcam and a marker card. When you hold the card up to the webcam, Flash is able to detect the orientation of the marker and superimpose, in this case, a 3d model on top of it. This technology has been around for a while now and having access to it in Flash is really exciting. Although it is still in its infancy, over time the speed and power of AR will improve to hopefully play an important role in how we view the web.

\section {Perancangan Aplikasi}
\label{sec:perancangan}
%%==========================================================================================================
% MULAI BAB IV
%==========================================================================================================
\chapter{PENGUJIAN DAN ANALISA}
\label{chap:pengujian}
%==========================================================================================================
% Subbab
%==========================================================================================================
Dalam pengujian aplikasi ini, ada dua cara pengujian yaitu dengan PC notebook yang dilengkapi \textit{webcam} internal atau menggunakan PC \textit{Desktop} dengan tambahan \textit{webcam}, konfigurasi sistem dapat dilihat pada gambar \ref{fig:monas_ar}. Untuk menjalankan aplikasi ini, diperlukan \textit{web browser} yang sudah terintegrasi \textit{plug-in} flash player versi 10. 

Spesifikasi sistem yang penulis gunakan dalam pengujian ini adalah sebagai berikut.
\begin{itemize}
\item Prosesor AMD Phenom\texttrademark II X4 955 Processor 3200 Mhz
\item Memori 2 Gb
\item standar usb 2.0 \textit{webcam} 
\item \textit{web browser} Mozilla Firefox versi 3.6.12
\end{itemize}

\begin{figure}[h]
\begin{center}
\includegraphics[width=6cm]{./images/ar3dv/monas_ar}
\caption{\label{fig:monas_ar} Hasil pengujian}
\end{center}
\end{figure}

Hasil pengujian dapat dilihat pada gambar \ref{fig:monas_ar}. Objek 3D dapat dirender dengan baik, meskipun kadang-kadang hilang dari marker, dikarenakan sistem tidak dapat menerima input dengan pergerakan yang marker cepat. Ada dua hal yang akan dianalisa yaitu fps (frame per second) dari video hasil rendering objek 3D dengan model yang berbeda-beda serta juga jarak dan sudut kemiringan marker yang masih dapat diterima aplikasi.

\section{Analisa \textit{Frame per Second (FPS)}}
\label{sec:analisa_fps}
\textit{Frame rate} atau frekuensi \textit{frame} merupakan frekuensi sebuah alat atau layar menghasilkan gambar yang disebut \textit{frame}. \textit{Frame rate} lebih sering dikenal dan diekspresikan sebagai frame per second (FPS). FPS merupakan jumlah gambar yang dihasilkan dalam waktu satu detik, mata manusia memerlukan jumlah \textit{frame} tertentu dalam satu detik agar dapat melihat gambar bergerak tanpa adanya \textit{flicker}. Jumlah FPS juga tergantung dari intensitas cahaya dan kecepatan objek bergerak dari video yang dihasilkan.

Pengujian untuk menganalisa FPS sangat diperlukan untuk mengetahui sejauh mana performa aplikasi viewer objek 3D ini. Untuk keperluan tersebut, penulis menambahkan informasi FPS tersebut ke aplikasi agar dapat dianalisa. Papervision3D telah menyediakan kelas untuk hal tersebut yaitu \textbf{StatsView}. Penggunaannya dapat dilihat dari \textit{listing} \ref{code:stats-view}.

\begin{lstlisting}[language=ActionScript,caption=StatsView,label=code:stats-view]
	_renderEngine = new LazyRenderEngine(_scene3D, _camera3D, _viewport3D);
	addChild(new StatsView(_renderEngine));
\end{lstlisting}

Sebuah informasi teks akan ditampilkan dikiri atas tampilan video, seperti yang ditunjukkan oleh gambar \ref{fig:fps_ar3dv}. Penulis menggunakan lima macam objek 3D berbeda dalam melakukan pengujian ini. Hasil pengujian dapat dilihat pada tabel \ref{tab:tabel_pengujian_fps}.

\begin{figure}[h]
\begin{center}
\includegraphics[width=6cm]{./images/ar3dv/monas_stats}
\caption{\label{fig:monas_stats} Hasil pengujian dengan statistik}
\end{center}
\end{figure}

\begin{table}[h]
\begin{center}
	\caption{Hasil Pengujian dengan Lima Model 3D}\label{tab:tabel_pengujian_fps}
	\vspace{5mm}
    \begin{tabular}[caption]{ | l | l | l | l | l | l | l |}
    \hline
    No & Nama & Tipe file & Ukuran file (KB) & FPS & Memory (MB) & Tri \cr \hline
    1 & airport & 3DS &  129,7 & 14 & 51,76 & 2870 \cr \hline
    2 & eiffel & DAE & 38,7 & 50 & 44,93 & 206 \cr \hline
	3 & monas & DAE & 154,5 & 24 & 53.80 & 564 \cr \hline
	4 & scout & DAE & 267,2 & 38 & 55.28 & 233 \cr \hline
	5 & tajmahal & DAE & 812,9 & 0 & 567.68 & 147.018 \cr \hline
    \hline
    \end{tabular}
\end{center}
\end{table}

Nilai FPS bergantung pada kompleksitas objek 3D yang ditampilkan. Untuk objek 3D yang tidak terlalu kompleks, hasilnya sangat baik, FPS berkisar antara 24 sampai 50, sehingga masih dapat terlihat baik. pada objek 3D kelima, nilai FPS nol atau tidak bergerak sama sekali, walaupun objek tersebut dapat terlihat. Penulis beranggapan bahwa frame yang dihasilkan dalam beberapa detik sekali, sehingga nilai FPS yang didapat adalah nol, hal ini disebabkan karena objek 3D yang di-\textit{render} sangat kompleks sehingga aplikasi tidak sanggup untum menampilkannya dengan baik.

%\section{Analisa Penggunaan Marker}
%\label{sec:analisa_marker}
%
%\subsection{Jarak}
%\label{subsec:analisa_jarak_marker}
%
%
%\subsection{Sudut Kemiringan}
%\label{subsec:analisa_sudut_marker}
%%==========================================================================================================
% MULAI BAB V
%==========================================================================================================
\chapter{KESIMPULAN DAN SARAN}
\label{chap:penutup}
%==========================================================================================================
% Subbab
%==========================================================================================================
\section{Kesimpulan}
\label{sec:kesimpulan}
Berdasarkan pembahasan bab-bab sebelumnya dan didukung oleh hasil pengujian, dapat diambil kesimpulan sebagai berikut.
\begin{enumerate}
\item Aplikasi dapat berjalan dengan baik tanpa perlu menginstall aplikasinya, karena aplikasi diakses menggunakan browser yang mempunyai \textit{plug-in} flash. 
\item Aplikasi dapat merender objek 3D dengan format DAE dan 3DS sesuai dengan marker yang didteksi, akan tetapi tingkat pengenalan marker sangat dipengaruhi cahaya dan bentuk marker yang digunakan.
\item Aplikasi dapat merender objek 3D sederhana dengan baik dengan nilai FPS yang relatif baik (24 sampai 50). 
\item Untuk objek 3D kompleks aplikasi tidak dapat berjalan dengan baik, terjadi \textit{flicker} atau gambar tidak bergerak sama sekali.
\end{enumerate}

\section{Saran}
\label{sec:saran}
\begin{enumerate}
\item Untuk pengembangan selanjutnya, diharapkan dapat menggunakan \textit{library} yang lebih baik. Marker yang digunakan diharapkan lebih variatif, dengan pola yang lebih menarik atau mungkin dengan pola yang berwarna sehingga lebih menarik.
\item Perangkat keras yang digunakan tidak lagi komputer tapi berupa \textit{gadget} yang lebih kecil dan bisa menjalankan flash, sehingga aplikasi benar-benar dapat dijalankan dari mana saja. 
\end{enumerate}

%==========================================================================================================
% Bagian daftar pustaka 
%==========================================================================================================
%\nocite{*}
\renewcommand{\bibname}{Daftar Pustaka}
\bibliographystyle{ieeetr} 
\bibliography{daftarpustaka}
\addcontentsline{toc}{chapter}{Daftar Pustaka} 

\end{document}

